\documentclass[a4paper,12pt]{article}

% 基础宏包配置
\usepackage{ctex}      % 中文支持
\usepackage{geometry}  % 页面布局
\usepackage{titlesec}  % 标题格式
\usepackage{fancyhdr}  % 页眉页脚
\usepackage{xcolor}    % 颜色支持
\usepackage{setspace}  % 行间距

% 页面边距设置 - 保持标准商务文档边距
\geometry{left=2.5cm, right=2.5cm, top=2.5cm, bottom=2.5cm}
\onehalfspacing % 1.5倍行距，增加文本块的厚重感

% 颜色定义
\definecolor{headerBlue}{RGB}{0, 51, 102}

% 标题格式定制
\titleformat{\section}
{\normalfont\Large\bfseries\color{headerBlue}}{\thesection}{1em}{}

% 页眉页脚
\pagestyle{fancy}
\fancyhf{}
\lhead{\small \textbf{2024年度ESG报告}}
\rhead{\small 管理层讨论与分析}
\cfoot{\thepage}

\title{\textbf{\huge 2024年度环境、社会及治理（ESG）报告}\\ \vspace{0.5em} \large 可持续发展管理与绩效综述}
\author{本公司董事会}
\date{2025年3月}

\begin{document}

\maketitle

\section{治理篇：稳健经营与合规管控}

2024年度，本公司在“三直四化”（电动化、智能网联化、高端化、国际化）战略指引下，积极应对全球宏观经济环境的不确定性，坚持高质量发展路径，实现了经营业绩的稳健增长与治理体系的持续优化。在经济绩效方面，受益于海外高端市场的持续拓展及国内市场需求的结构性复苏，公司全年实现营业收入372.18亿元，较上年同期增长37.63\%，归属于母公司股东的净利润达到41.16亿元。公司董事会始终坚持为股东创造可持续价值，高度重视股东回报机制的稳定性与连续性，报告期内公司实施了两次现金分红方案，全年累计派发现金股利总额达44.28亿元，分红总额超过当期归母净利润，充分体现了公司良好的现金流状况及对投资者权益的切实保障。

在合规管理与内部控制方面，公司持续深化合规体系建设，确保企业运营严格遵循法律法规及监管要求。报告期内，公司凭借规范的公司治理结构与高质量的信息披露表现，连续第13年获得证券交易所信息披露工作“A级”评价，并顺利通过了ISO 37301合规管理体系认证，标志着公司合规管理架构已全面接轨国际标准。供应链管理作为公司ESG治理的重要组成部分，公司致力于构建绿色、廉洁、具有韧性的供应链生态。截至2024年12月31日，公司合格供应商共计530家，其中河南本地供应商占比约为20\%，有效发挥了龙头企业对区域产业链的带动作用。为强化商业道德风险管控，公司严格执行廉洁准入机制，报告期内已与超过900家供应商（含潜在供应商）签署了廉洁协议，明确了双方在反商业贿赂及反腐败方面的权利义务，从源头营造风清气正的商业环境。

\section{社会篇：研发创新与利益相关方责任}

公司坚持将技术创新作为驱动可持续发展的核心引擎，持续加大在新能源及智能网联领域的资源投入。2024年度，公司研发投入金额为17.88亿元，占营业收入的比例为4.80\%。公司拥有一支高素质的研发队伍，研发人员总数为3,505人，其中博士学历人员16人，为技术突破提供了坚实的人才支撑。在知识产权与标准建设方面，报告期内公司新增授权专利176项，其中发明专利114项，累计参与制定国家及行业标准317项，持续引领行业技术规范的发展。在核心技术领域，公司基于“软件定义汽车”理念，自主研发了YOS车机操作系统，实现了软硬件解耦与全栈OTA升级能力；全面应用了高集成度的C架构（中央计算+区域控制），有效提升了电子电气架构的通信效率与轻量化水平。在电驱动技术方面，公司开发的多合一动力域控制器采用了碳化硅（SiC）功率模块，相较于传统硅基模块，其开关损耗降低了50\%至70\%，显著提升了系统综合能效。针对新能源车辆冬季续航衰减的行业痛点，公司推出了多源超低温热泵系统，通过高效回收环境热能及电机余热，在低温工况下采暖节能效果达到30\%。此外，在智能网联应用端，公司L4级自动驾驶巴士已累计安全运营超过1,700万公里，并成功获批国家首批智能网联汽车准入试点，验证了自动驾驶技术的安全性与商业化潜力。

在人力资源管理方面，截至报告期末，公司员工总数为16,707人，其中女性员工占比为10.02\%。公司严格遵守《劳动法》等法律法规，建立健全职业健康安全管理体系。报告期内，公司在安全生产领域的投入持续增加，其中劳动防护用品投入5,500余万元，生产现场环境改善投入2,200余万元，安全教育培训投入180余万元。通过全年开展7.9万人次的安全培训及隐患排查治理，公司实现了全年无新增职业病的安全绩效目标。在履行企业公民责任方面，公司2024年度对外捐赠及公益项目总投入超过4,000万元。公司持续开展“儿童交通安全公益行”项目，足迹覆盖全国13个省市，通过车辆盲区体验等实景教学活动，惠及近万名儿童。在海外市场，公司结合业务布局在南美某国实施了“零碳森林”项目，累计种植耐旱树种36,000余棵，在助力当地生态修复的同时，积极响应全球应对气候变化的行动倡议。

\section{环境篇：绿色运营与资源效率}

公司承诺在运营全过程中贯彻绿色低碳理念，严格管控能源消耗与污染物排放，致力于降低产品全生命周期的环境足迹。2024年度，公司建立了覆盖主要生产基地的碳排放数据监测与核算体系。根据核算结果，报告期内公司范围1（直接温室气体排放）为57,330.53吨二氧化碳当量，范围2（能源间接温室气体排放）为122,786.63吨二氧化碳当量，温室气体排放强度为0.07吨二氧化碳/万元营收。公司综合能源消费量为34,560.94吨标准煤。

为实现节能降碳目标，公司在报告期内大力推进节能技术改造与清洁能源替代。生产制造环节，公司全年实施重点节能改造项目7项，通过设备能效升级与工艺优化，实现年节电1,258兆瓦时，年节气11.34万立方米；同时，公司完善了工业余热回收系统，全年回收利用余热约29,306吉焦，有效减少了化石能源消耗。在清洁能源利用方面，公司新能源智造基地屋顶光伏发电系统全年发电量达155.68万千瓦时，折合减少碳排放约835吨；此外，公司全年通过电力市场采购绿色电力13,580兆瓦时，进一步优化了能源消费结构。

在资源循环利用与废弃物管理方面，公司坚持减量化、资源化、无害化原则。水资源管理上，公司2024年新取水量为141.25万吨，通过完善的中水回用设施，全年重复用水量达到6,465.36万吨，重复用水率高达97.8\%，基本实现了工业用水的内部循环。废弃物处置方面，公司严格执行分类管理，报告期内一般工业固废中回收利用量为61,485.76吨，危险废物产生量为4,110.84吨，全部委托具备相应资质的第三方机构进行无害化处置，合规处置率保持100\%。绿色物流方面，公司在零部件采购与运输环节大力推广可循环包装应用，目前国产零部件周转包装占比已达到89\%，显著降低了供应链物流环节的一次性包装废弃物产生量。展望未来，公司将继续探索氢能技术应用及供应链碳足迹管理，推动实现更深层次的绿色转型。

\vspace{3cm}

\end{document}